\documentclass[12pt,a4paper]{article}
\usepackage[UTF8]{ctex}
\usepackage{amsmath,amssymb,amsthm}
\usepackage{geometry}
\usepackage{graphicx}

\geometry{margin=2.5cm}

\title{结合轨迹跟踪的导纳控制的稳态误差分析}
\author{第11章补充说明}
\date{}

\begin{document}

\maketitle

\section{问题提出}

结合轨迹跟踪的导纳控制公式为：
\begin{equation}
M\ddot{x}_d + B(\dot{x}_d - \dot{x}) + K(x_d - x) = f_{\text{ext}} \tag{*}
\end{equation}

\textbf{关键问题}：这个公式对于位置精度来说存在稳态误差。本文将详细分析稳态误差的原因、大小和解决方法。

\section{稳态误差分析}

\subsection{稳态条件}

在稳态时，假设：
\begin{itemize}
\item 期望位置和速度恒定：$\dot{x}_d = 0$，$\ddot{x}_d = 0$
\item 实际速度为零：$\dot{x} = 0$
\item 外力恒定：$f_{\text{ext}} = f_0$（常数）
\end{itemize}

\subsection{稳态误差方程}

在稳态条件下，公式(*)简化为：
\begin{equation}
B \cdot 0 + K(x_d - x_{ss}) = f_0
\end{equation}

即：
\begin{equation}
K x_{e,ss} = f_0
\end{equation}

其中$x_{e,ss} = x_d - x_{ss}$是稳态位置误差。

\textbf{稳态误差}：
\begin{equation}
x_{e,ss} = \frac{f_0}{K} \tag{1}
\end{equation}

\subsection{关键结论}

\begin{enumerate}
\item \textbf{存在稳态误差}：
\begin{itemize}
\item 如果$f_{\text{ext}} \neq 0$，则$x_{e,ss} \neq 0$
\item 稳态误差与外力成正比
\item 稳态误差与刚度$K$成反比
\end{itemize}

\item \textbf{误差大小}：
\begin{itemize}
\item 外力越大，误差越大
\item 刚度越大，误差越小
\item 但增大刚度可能导致系统"硬"，对外力响应不灵敏
\end{itemize}
\end{enumerate}

\section{物理意义}

\subsection{力的平衡}

公式(*)可以理解为力的平衡：
\begin{equation}
\text{虚拟弹簧力} + \text{虚拟阻尼力} + \text{虚拟惯性力} = \text{外力}
\end{equation}

在稳态时：
\begin{equation}
K(x_d - x) = f_{\text{ext}}
\end{equation}

\textbf{物理意义}：
\begin{itemize}
\item 虚拟弹簧试图将机器人拉向期望位置$x_d$
\item 但外力$f_{\text{ext}}$阻止这个运动
\item 在平衡时，弹簧力等于外力
\item 因此存在位置误差$x_e = f_{\text{ext}}/K$
\end{itemize}

\subsection{类比：弹簧-质量系统}

想象一个弹簧-质量系统：
\begin{itemize}
\item 弹簧一端固定在期望位置$x_d$
\item 质量在位置$x$
\item 外力$f_{\text{ext}}$作用在质量上
\end{itemize}

在平衡时：
\begin{equation}
k(x_d - x) = f_{\text{ext}} \quad \Rightarrow \quad x_d - x = \frac{f_{\text{ext}}}{k}
\end{equation}

这就是稳态误差的来源。

\section{稳态误差的影响因素}

\subsection{外力$f_{\text{ext}}$的影响}

从公式(1)：
\begin{equation}
x_{e,ss} = \frac{f_{\text{ext}}}{K}
\end{equation}

\textbf{结论}：
\begin{itemize}
\item 外力越大，稳态误差越大
\item 如果$f_{\text{ext}} = 0$，则$x_{e,ss} = 0$（无误差）
\item 这是导纳控制的特性：允许外力影响位置
\end{itemize}

\subsection{刚度$K$的影响}

\textbf{高刚度$K$}：
\begin{itemize}
\item 稳态误差小：$x_{e,ss} = f_{\text{ext}}/K$（小）
\item 但系统"硬"，对外力响应不灵敏
\item 接近纯轨迹跟踪控制
\end{itemize}

\textbf{低刚度$K$}：
\begin{itemize}
\item 稳态误差大：$x_{e,ss} = f_{\text{ext}}/K$（大）
\item 但系统"软"，对外力响应灵敏
\item 接近标准导纳控制
\end{itemize}

\textbf{权衡}：
\begin{itemize}
\item 需要在位置精度和力响应之间权衡
\item 不能同时获得高精度和高灵敏度
\end{itemize}

\section{与标准导纳控制的对比}

\subsection{标准导纳控制}

标准导纳控制公式：
\begin{equation}
M\ddot{x}_d + B\dot{x} + Kx = f_{\text{ext}}
\end{equation}

\textbf{特点}：
\begin{itemize}
\item 没有期望位置$x_d$
\item 不存在"位置误差"的概念
\item 位置由外力和初始条件决定
\item 可能漂移
\end{itemize}

\subsection{结合轨迹跟踪的导纳控制}

公式：
\begin{equation}
M\ddot{x}_d + B(\dot{x}_d - \dot{x}) + K(x_d - x) = f_{\text{ext}}
\end{equation}

\textbf{特点}：
\begin{itemize}
\item 有期望位置$x_d$
\item 存在位置误差$x_e = x_d - x$
\item 在稳态时，误差$x_{e,ss} = f_{\text{ext}}/K$
\item 位置不会漂移（有期望位置锚定）
\end{itemize}

\subsection{对比总结}

\begin{table}[h]
\centering
\begin{tabular}{|l|l|l|}
\hline
\textbf{方面} & \textbf{标准导纳控制} & \textbf{结合轨迹跟踪} \\
\hline
期望位置 & 无 & 有$x_d$ \\
\hline
位置误差 & 无概念 & 有$x_e = x_d - x$ \\
\hline
稳态误差 & 无概念 & $x_{e,ss} = f_{\text{ext}}/K$ \\
\hline
位置漂移 & 可能 & 不会 \\
\hline
\end{tabular}
\caption{两种导纳控制的对比}
\end{table}

\section{减少稳态误差的方法}

\subsection{方法1：增大刚度$K$}

\textbf{优点}：
\begin{itemize}
\item 直接减少稳态误差：$x_{e,ss} = f_{\text{ext}}/K$
\item 提高位置精度
\end{itemize}

\textbf{缺点}：
\begin{itemize}
\item 系统变"硬"，对外力响应不灵敏
\item 可能影响人机交互体验
\item 可能导致系统不稳定（如果$K$太大）
\end{itemize}

\subsection{方法2：使用积分项}

在控制律中加入积分项，类似于PID控制：

\textbf{修改后的公式}：
\begin{equation}
M\ddot{x}_d + B(\dot{x}_d - \dot{x}) + K(x_d - x) + K_i \int_0^t (x_d - x) d\tau = f_{\text{ext}} \tag{**}
\end{equation}

\textbf{稳态分析}：

在稳态时，如果$f_{\text{ext}}$是常数：
\begin{equation}
K x_{e,ss} + K_i \int_0^t x_{e,ss} d\tau = f_0
\end{equation}

对时间求导：
\begin{equation}
K \dot{x}_{e,ss} + K_i x_{e,ss} = 0
\end{equation}

如果系统稳定，$\dot{x}_{e,ss} = 0$，因此：
\begin{equation}
K_i x_{e,ss} = 0
\end{equation}

由于$K_i > 0$，必须有$x_{e,ss} = 0$。

\textbf{结论}：积分项可以消除稳态误差！

\textbf{但需要注意}：
\begin{itemize}
\item 积分项可能影响系统稳定性
\item 需要适当选择$K_i$
\item 可能需要积分器抗饱和
\end{itemize}

\subsection{方法3：前馈补偿}

如果能够估计或测量外力$f_{\text{ext}}$，可以使用前馈补偿：

\textbf{修改后的公式}：
\begin{equation}
M\ddot{x}_d + B(\dot{x}_d - \dot{x}) + K(x_d - x) = f_{\text{ext}} - \hat{f}_{\text{ext}}
\end{equation}

其中$\hat{f}_{\text{ext}}$是外力的估计值。

\textbf{稳态分析}：

如果$\hat{f}_{\text{ext}} = f_{\text{ext}}$（完美估计），则：
\begin{equation}
K x_{e,ss} = 0 \quad \Rightarrow \quad x_{e,ss} = 0
\end{equation}

\textbf{优点}：
\begin{itemize}
\item 可以消除稳态误差
\item 不需要积分项
\end{itemize}

\textbf{缺点}：
\begin{itemize}
\item 需要准确估计外力
\item 可能难以实现
\end{itemize}

\subsection{方法4：自适应刚度}

根据外力大小自适应调整刚度：

\begin{equation}
K = \begin{cases}
K_{\text{high}} & \text{如果} |f_{\text{ext}}| < f_{\text{threshold}} \\
K_{\text{low}} & \text{如果} |f_{\text{ext}}| \geq f_{\text{threshold}}
\end{cases}
\end{equation}

\textbf{思路}：
\begin{itemize}
\item 小外力时：高刚度，高精度
\item 大外力时：低刚度，允许偏离
\end{itemize}

\section{数值例子}

\subsection{例子1：恒定外力}

\textbf{参数}：
\begin{itemize}
\item 期望位置：$x_d = 0.5$ m
\item 外力：$f_{\text{ext}} = 10$ N
\item 刚度：$K = 100$ N/m
\end{itemize}

\textbf{稳态误差}：
\begin{equation}
x_{e,ss} = \frac{10}{100} = 0.1 \text{ m}
\end{equation}

\textbf{实际位置}：
\begin{equation}
x_{ss} = x_d - x_{e,ss} = 0.5 - 0.1 = 0.4 \text{ m}
\end{equation}

\textbf{结论}：机器人停在$0.4$ m处，而不是期望的$0.5$ m。

\subsection{例子2：增大刚度}

如果增大刚度到$K = 1000$ N/m：

\textbf{稳态误差}：
\begin{equation}
x_{e,ss} = \frac{10}{1000} = 0.01 \text{ m} = 1 \text{ cm}
\end{equation}

\textbf{实际位置}：
\begin{equation}
x_{ss} = 0.5 - 0.01 = 0.49 \text{ m}
\end{equation}

\textbf{结论}：误差从$10$ cm减少到$1$ cm，但系统变"硬"。

\subsection{例子3：使用积分项}

如果加入积分项$K_i = 10$ N/(m·s)：

\textbf{稳态误差}：
\begin{equation}
x_{e,ss} = 0 \text{ m}
\end{equation}

\textbf{实际位置}：
\begin{equation}
x_{ss} = x_d = 0.5 \text{ m}
\end{equation}

\textbf{结论}：稳态误差为零，但需要保证系统稳定。

\section{实际应用考虑}

\subsection{何时可以接受稳态误差？}

\begin{enumerate}
\item \textbf{人机协作任务}：
\begin{itemize}
\item 允许操作员推拉调整位置
\item 稳态误差是"特性"而不是"缺陷"
\item 误差大小可以通过刚度调整
\end{itemize}

\item \textbf{装配任务}：
\begin{itemize}
\item 如果外力是预期的（如装配阻力）
\item 稳态误差可能可以接受
\item 但如果需要高精度，需要消除误差
\end{itemize}
\end{enumerate}

\subsection{何时需要消除稳态误差？}

\begin{enumerate}
\item \textbf{高精度定位}：
\begin{itemize}
\item 需要精确到达目标位置
\item 必须消除稳态误差
\item 可以使用积分项或前馈补偿
\end{itemize}

\item \textbf{轨迹跟踪}：
\begin{itemize}
\item 如果外力是干扰（非预期）
\item 应该消除其影响
\item 使用积分项或增大刚度
\end{itemize}
\end{enumerate}

\section{稳定性考虑}

\subsection{使用积分项时的稳定性}

加入积分项后，系统变为三阶：
\begin{equation}
M\ddot{x}_d + B(\dot{x}_d - \dot{x}) + K(x_d - x) + K_i \int_0^t (x_d - x) d\tau = f_{\text{ext}}
\end{equation}

\textbf{稳定性条件}：
\begin{itemize}
\item 需要适当选择$K_i$
\item $K_i$太大可能导致不稳定
\item 可以使用Routh-Hurwitz判据分析
\end{itemize}

\subsection{增大刚度时的稳定性}

\textbf{问题}：
\begin{itemize}
\item 高刚度可能导致系统"硬"
\item 在接触任务中可能不稳定
\item 需要适当的阻尼$B$
\end{itemize}

\section{总结}

\subsection{核心结论}

\begin{enumerate}
\item \textbf{存在稳态误差}：
\begin{equation}
x_{e,ss} = \frac{f_{\text{ext}}}{K}
\end{equation}

\item \textbf{误差特性}：
\begin{itemize}
\item 与外力成正比
\item 与刚度成反比
\item 这是导纳控制的固有特性
\end{itemize}

\item \textbf{减少误差的方法}：
\begin{itemize}
\item 增大刚度$K$（但系统变硬）
\item 使用积分项（但可能影响稳定性）
\item 前馈补偿（但需要估计外力）
\item 自适应刚度（但实现复杂）
\end{itemize}
\end{enumerate}

\subsection{设计权衡}

\begin{itemize}
\item \textbf{位置精度 vs. 力响应灵敏度}：
\begin{itemize}
\item 高精度需要高刚度，但响应不灵敏
\item 高灵敏度需要低刚度，但精度低
\end{itemize}

\item \textbf{稳态误差 vs. 系统稳定性}：
\begin{itemize}
\item 使用积分项可以消除误差，但可能不稳定
\item 需要仔细设计参数
\end{itemize}

\item \textbf{实际应用需求}：
\begin{itemize}
\item 根据应用需求决定是否可以接受稳态误差
\item 人机协作可能可以接受，高精度定位不能接受
\end{itemize}
\end{enumerate}

\subsection{关键公式}

\textbf{稳态误差}：
\begin{equation}
x_{e,ss} = \frac{f_{\text{ext}}}{K}
\end{equation}

\textbf{带积分项的控制}：
\begin{equation}
M\ddot{x}_d + B(\dot{x}_d - \dot{x}) + K(x_d - x) + K_i \int_0^t (x_d - x) d\tau = f_{\text{ext}}
\end{equation}

\textbf{稳态时}（如果稳定）：$x_{e,ss} = 0$

\end{document}

